% ============================================================
% Come si usano, argomento 17: processi aleatori
% ============================================================
\argomento{Processi aleatori}{Formule a p.~\pageref*{f:catenavettori}}

\begin{scelta}{Come scegliere la formula}
\renewcommand{\arraystretch}{1.25}
\begin{tabular}{@{}P{0.55\linewidth}P{0.39\linewidth}@{}}
  \toprule
  \textbf{Se il testo\ldots} & \textbf{si usa\ldots} \\
  \midrule
  chiede la probabilità di una sequenza di valori, uno dopo l'altro & la regola della catena \vv{catenavettori};
    con campioni indipendenti il prodotto \vv{procindip} \\
  chiede se un processo è stazionario (del primo ordine, del secondo, in senso stretto) & \vv{stazprimo}
    \vv{stazsecondo} \vv{stazstretta} \\
  chiede se un processo è SSL, o la sua potenza media & \vv{autocorrelazione} \vv{ssl} \\
  chiede la matrice di covarianza di alcuni campioni & \vv{toeplitz} \\
  parla di bit $\pm 1$ indipendenti, di medie di campioni vicini, di un oscillatore, di una pedina che va avanti e
    indietro & i processi notevoli \vv{bernoulliproc} \vv{mediamobile} \vv{sinusoideproc} \vv{passeggiata} \\
  \bottomrule
\end{tabular}

Per decidere se un processo è SSL si calcolano $E[X(n)]$ e $E[X(n)\, X(n + k)]$: se $n$ non compare in nessuna delle
due, il processo è SSL. Per dire di no basta un controesempio: due coppie di istanti alla stessa distanza con
autocorrelazioni diverse.
\end{scelta}

\scheda{catenavettori}{
  \usoquando{Le variabili «escono» una dopo l'altra (estrazioni senza reinserimento, campioni di un segnale) e ogni
    condizionata è facile da scrivere.}
  \usocome{Si moltiplicano la probabilità della prima, quella della seconda sapendo la prima, quella della terza
    sapendo le prime due, e così via.}
  \usoesempio{Tre assi di fila da un mazzo di 52 carte: $\frac{4}{52} \cdot \frac{3}{51} \cdot \frac{2}{50} =
    \frac{24}{132\,600} \simeq 1{,}8 \cdot 10^{-4}$.}
}

\scheda{processo}{
  \usoquando{Il risultato dell'esperimento è un intero segnale nel tempo: un rumore, una sequenza di bit, una misura
    ripetuta ogni giorno.}
  \usocome{Lo si guarda in due modi. Fissato il risultato, è una successione di numeri (quella che si vede
    sull'oscilloscopio); fissato l'istante, è una variabile aleatoria, con la sua densità.}
  \usoesempio{Sinusoide con fase a caso: ogni realizzazione è una sinusoide perfetta; il valore all'istante $n_0$, al
    variare della fase, è una variabile aleatoria fra $-A$ e $A$.}
}

\scheda{stazprimo}{
  \usoquando{Si chiede se i campioni hanno tutti la stessa distribuzione, istante per istante.}
  \usocome{Si calcola la densità (o la pmf) di $X(n)$ con $n$ generico e si guarda se $n$ compare.}
  \usoesempio{Sinusoide con fase a caso: per ogni $n$ la densità è l'arcoseno $\frac{1}{\pi\sqrt{A^2 - x^2}}$, che
    non dipende da $n$. La passeggiata aleatoria invece non lo è: la pmf di $S(n)$ si allarga con $n$.}
}

\scheda{stazsecondo}{
  \usoquando{Si vuole sapere se i legami fra due campioni dipendono solo dalla loro distanza.}
  \usocome{Si scrive la congiunta di $X(n_1)$ e $X(n_2)$ e si controlla che dipenda solo da $n_2 - n_1$. Se il
    processo è stazionario del secondo ordine lo è anche del primo: la marginale si ottiene integrando la congiunta.}
  \usoattenzione{Il primo ordine non basta. Il processo di Bernoulli e quello che ripete per sempre il primo lancio,
    $X(n) = X(0)$, hanno la stessa pmf del primo ordine ($\pm 1$ con probabilità $\frac{1}{2}$), ma nel secondo i
    campioni sono tutti uguali.}
  \usoesempio{Processo di Bernoulli: $P\big(X(n_1) = x_1, X(n_2) = x_2\big) = \frac{1}{4}$ per ogni coppia di valori
    e di istanti distinti.}
}

\scheda{stazstretta}{
  \usoquando{Si chiede la stazionarietà «completa»: nessuna statistica del processo cambia se si sposta l'origine dei
    tempi.}
  \usocome{Per un processo indipendente basta controllare che la densità (o la pmf) di un campione sia la stessa in
    ogni istante (\vv{procindip}).}
  \usoesempio{Il processo di Bernoulli e il quaternario indipendente sono stazionari in senso stretto.}
}

\scheda{procindip}{
  \usoquando{I campioni sono indipendenti (bit trasmessi, lanci ripetuti): la caratterizzazione completa si riduce a
    quella del primo ordine.}
  \usocome{Si moltiplicano le pmf (o le densità) dei singoli campioni; se il processo è anche stazionario, sono tutte
    la stessa $p_X$.}
  \usoesempio{Bernoulli $\pm 1$: la sequenza $+1, +1, -1$ in tre istanti qualsiasi ha probabilità $\left(
    \frac{1}{2} \right)^3 = \frac{1}{8}$. Quaternario: due campioni entrambi uguali a 2 hanno probabilità
    $\frac{1}{4} \cdot \frac{1}{4} = \frac{1}{16}$.}
}

\scheda{autocorrelazione}{
  \usoquando{Serve una caratterizzazione sintetica del processo: quanto si somigliano, in media, i campioni in due
    istanti.}
  \usocome{Si scrive il prodotto $X(n_1)\, X(n_2)$, lo si sviluppa e si usa la linearità della media.}
  \usoesempio{Bernoulli $\pm 1$: con $n_1 = n_2$ il prodotto è $X^2(n) = 1$, quindi $R_X(n, n) = 1$; con
    $n_1 \neq n_2$ i campioni sono indipendenti e $R_X(n_1, n_2) = E[X(n_1)]\, E[X(n_2)] = 0$.}
}

\scheda{ssl}{
  \usoquando{Si chiede se il processo è stazionario in senso lato, oppure servono la potenza media o la correlazione
    fra campioni a distanza $k$.}
  \usocome{Si calcolano $E[X(n)]$ e $E[X(n)\, X(n + k)]$: il processo è SSL se la prima non dipende da $n$ e la
    seconda dipende solo da $k$.}
  \usoattenzione{SSL non vuol dire stazionario. Campioni indipendenti $\mathcal{N}(0, 1)$ negli istanti pari e
    $\mathcal{U}(-\sqrt{3}, \sqrt{3})$ in quelli dispari hanno media 0, varianza 1 e $R_X(n_1, n_2) = 0$ per
    $n_1 \neq n_2$: è SSL, ma la densità cambia da un istante all'altro.}
  \usoesempio{Passeggiata aleatoria: la media è costante, ma $R_S(1, 2) = 1$ e $R_S(10, 11) = 10$ con la stessa
    distanza: non è SSL.}
}

\scheda{toeplitz}{
  \usoquando{Si prelevano alcuni campioni da un processo SSL e servono il vettore delle medie e la matrice di
    covarianza.}
  \usocome{Sulla diagonale va la varianza $\sigma_X^2 = R_X(0) - \mu^2$, nella posizione $(i, j)$ la covarianza
    $R_X(n_j - n_i) - \mu^2$.}
  \usoesempio{Media mobile, quattro campioni consecutivi ($\sigma_Y^2 = \frac{1}{2}$, covarianza $\frac{1}{4}$ fra
    campioni vicini): $\matr{C}_{\vett{Y}} = \frac{1}{4} \begin{bmatrix} 2 & 1 & 0 & 0 \\ 1 & 2 & 1 & 0 \\ 0 & 1 & 2
    & 1 \\ 0 & 0 & 1 & 2 \end{bmatrix}$: ogni diagonale ha elementi uguali.}
}

\scheda{bernoulliproc}{
  \usoquando{È il modello di una sequenza di bit trasmessi con i livelli $\pm 1$, e il mattone con cui si costruiscono
    altri processi (media mobile, passeggiata aleatoria).}
  \usocome{Media e autocorrelazione vengono da $X^2(n) = 1$ e dall'indipendenza dei campioni.}
  \usoesempio{Un campione non dice nulla sugli altri: la matrice di covarianza di qualsiasi vettore di campioni è
    l'identità. Il processo è stazionario in senso stretto (\vv{procindip}).}
}

\scheda{mediamobile}{
  \usoquando{Il processo si ottiene mediando campioni vicini di un processo indipendente: i campioni che hanno un
    ingresso in comune diventano correlati.}
  \usocome{Si sviluppa il prodotto $Y(n)\, Y(n + k)$ e si tengono solo i termini con lo stesso campione di $X$, che
    valgono 1.}
  \usoesempio{$Y(n)$ vale $-1$, 0, $+1$ con probabilità $\frac{1}{4}$, $\frac{1}{2}$, $\frac{1}{4}$. Con $k = 1$
    l'unico campione in comune fra $Y(n)$ e $Y(n + 1)$ è $X(n)$, quindi $R_Y(1) = \frac{1}{4}$.}
}

\scheda{sinusoideproc}{
  \usoquando{Un oscillatore che si accende in un istante qualsiasi: ogni realizzazione è una sinusoide, il caso decide
    solo la fase.}
  \usocome{La media è un coseno mediato su un periodo intero, quindi nulla; per l'autocorrelazione si usa il prodotto
    di due coseni (\vv{prodottocoseni}).}
  \usoesempio{Con $f_0 = \frac{1}{8}$, un periodo ogni 8 campioni: $R_X(4) = \frac{A^2}{2} \cos\pi = -\frac{A^2}{2}$,
    campioni a mezzo periodo opposti; $R_X(8) = \frac{A^2}{2}$, a un periodo di nuovo uguali. La memoria non si esaurisce
    mai.}
}

\scheda{passeggiata}{
  \usoquando{Una grandezza che si accumula passo dopo passo con incrementi indipendenti: una pedina che avanza o
    arretra, un errore che si somma a ogni misura.}
  \usocome{La media con la linearità, la varianza sommando quelle dei passi (\vv{varsomma}), l'autocorrelazione
    contando i termini $E[X(i)\, X(j)]$ con $i = j$.}
  \usoesempio{Dopo 100 passi la deviazione standard è $\sqrt{100} = 10$: la pedina è tipicamente a una decina di
    caselle dall'origine. $R_S(10, 11) = \min(10, 11) = 10$.}
}
